% Appendix D: Experiment 3 details. Design part written 2026-09-11 from
% code/experiments/exp3_real_context.py, exp3_train.py,
% results/exp3/readiness_scan.csv and results/exp3/calibration3.json.

\subsection{Data}
Breast-tissue QCL micro-array cores, fold~0 of the five-fold patient-level
split used throughout this work ($115$ training, $28$ validation, $26$ test
cores). Each labelled tissue pixel contributes the $942$-dimensional vector
of absorbance, first and second derivative ($3\times314$ channels), after
the field-standard preprocessing of the companion pipeline, which includes
per-core PCA denoising with $23$ components before the derivatives are
taken; each core's spectra are therefore of numerical rank $23$, while the
pooled training-patient data are not, since the retained subspaces differ
between cores. Up to
$600$ (training) or $300$ (validation, test) pixels per class per core were
sampled with a fixed seed: $75{,}297$ training pixels from $79$ cores,
$11{,}024$ validation from $20$, $8{,}038$ test from $16$, over the four
classes; the experiment uses the pair Cancer Epithelium ($+1$; $29{,}603$
training pixels) versus Cancer-Associated Stroma ($-1$; $21{,}923$).

\subsection{Readiness scan}
For each class pair, bottleneck $K$ and preprocessing, Table~\ref{tab:exp3_scan}
reports the centre-only oracle (ridge Fisher discriminant on the preprocessed
$942$-vector, fitted on one half of the training cores and evaluated on the
other half) and the same discriminant through three random encoders
$W\sim\mathcal N(0,1/942)^{K\times942}$. Standardization is per feature;
whitening is PCA whitening with a relative ridge of $10^{-3}$ on the
eigenvalues; both are estimated on the fitting half.

\begin{table}[h]
\centering\small
\caption{Readiness scan (\texttt{results/exp3/readiness\_scan.csv}): centre-only
oracle and random-encoder probe (mean over three encoders) by pair, $K$ and
preprocessing. $|\cos|$ is the cosine between the class-mean contrast and the
first principal direction of the standardized training half.}
\label{tab:exp3_scan}
\begin{tabular}{llcccccc}
\toprule
pair & preprocessing & oracle & $K=2$ & $K=4$ & $K=12$ & $K=32$ & $|\cos|$ \\
\midrule
CancerEpi vs CAS & standardized & 0.969 & 0.726 & 0.803 & 0.955 & 0.968 & 0.50 \\
CancerEpi vs CAS & whitened & 0.944 & 0.503 & 0.486 & 0.539 & 0.591 & 0.50 \\
NormalStroma vs CAS & standardized & 0.810 & 0.580 & 0.639 & 0.674 & 0.741 & 0.44 \\
NormalStroma vs CAS & whitened & 0.768 & 0.502 & 0.509 & 0.529 & 0.600 & 0.44 \\
NormalEpi vs CancerEpi & standardized & 0.758 & 0.547 & 0.573 & 0.662 & 0.724 & 0.20 \\
NormalEpi vs CancerEpi & whitened & 0.761 & 0.575 & 0.545 & 0.624 & 0.580 & 0.20 \\
NormalStroma vs CancerEpi & standardized & 0.989 & 0.781 & 0.874 & 0.978 & 0.990 & 0.88 \\
NormalStroma vs CancerEpi & whitened & 0.974 & 0.579 & 0.541 & 0.549 & 0.626 & 0.88 \\
\bottomrule
\end{tabular}
\end{table}

With standardized inputs a random encoder already reads the
epithelium-versus-stroma contrasts, because they lie along the top
principal directions; whitening removes that accessibility while the
oracle stays high. The two regimes of Section~\ref{sec:real} follow: the
same pair, $K=12$, standardized (\emph{ready}) or whitened (\emph{unready}).

\subsection{Patch construction}
Each training example is a real centre pixel with its recorded label and
eight real donor spectra sampled uniformly from the same split side (resampled
once if the donor's core equals the centre's), independent of the centre
label. In the preprocessed space, donor $d$ receives the constructed cue
$\gamma_d\,(c\,s+\tau\eta_d)\,e_d$, where $e_1,\dots,e_8$ are the first
eight principal coordinates (in the standardized regime the directions
$v_1,\dots,v_8$ with $\gamma_d=3\sqrt{\lambda_d}$; in the whitened regime the
first eight whitened coordinates with $\gamma_d=30$), $c$ is the centre's
label code, and $s=+1$ (informative), $s=-1$ (reversed) or $c$ replaced by an
independent random label (context-random); \emph{spectral-only} removes the
cue and keeps the donors. The draw order is condition-independent, so the
four conditions generated from one seed share centres, donors and $\eta$.
Preprocessing, principal directions and $\lambda_d$ are estimated on
training pixels only. Training: $24{,}000$ class-balanced centres; probe
fitting: $6{,}000$ disjoint training pixels; evaluation: up to $6{,}000$
validation and $6{,}000$ test centres with donors from their own side.
$\tau$ is calibrated per regime on training-core halves so that the
neighbourhood oracle (a discriminant on the eight cue projections) matches
the centre-only oracle.

\begin{table}[h]
\centering\small
\caption{Experiment 3 calibration (\texttt{results/exp3/calibration3.json}).
Readiness: discriminant through a random $12\times942$ encoder on the
centre alone and on the encoded $3\times3$ patch (three encoders).}
\label{tab:exp3_calib}
\begin{tabular}{lcccccc}
\toprule
regime & $\tau$ & centre oracle (fit/val) & context oracle & random enc.\ centre & random enc.\ patch \\
\midrule
ready (standardized) & 1.475 & 0.969 / 0.959 & 0.969 & 0.935, 0.970, 0.952 & 0.988, 0.994, 0.991 \\
unready (whitened) & 1.735 & 0.948 / 0.884 & 0.948 & 0.476, 0.512, 0.605 & 0.948, 0.947, 0.950 \\
\bottomrule
\end{tabular}
\end{table}

\subsection{Model, training and arms}
Encoder $W\in\R^{12\times942}$ applied to the nine spectra; head
$h=\mathrm{ReLU}(W_1\,\mathrm{vec}(z)+b_1)$ of width $M$ over the $9K=108$
encoded coordinates and a linear readout to one logit; margin loss;
full-batch gradient descent with $\eta=10^{-3}$, at most $40{,}000$ steps,
the same thresholds and snapshot rule as Experiment~2. Arms per regime:
\texttt{sp} at $M\in\{8,32,128,512\}$; \texttt{headlr} with
$\kappa\in\{1,1/16,1/256\}$ at $M=32$; \texttt{ctxfree} (context-random
training) at $M\in\{8,128\}$; \texttt{frozen} at $M=32$; seeds $0$--$2$.
Spectral probe: ridge discriminant on the encoder output of centre pixels,
fitted on the probe set and evaluated on validation (and test) centres,
reported as the gain over its initial value. Reliance and failure are
measured on validation patients with the paired conditions; test patients
are reported separately and were not used for any choice.

\subsection{Results}
Tables generated by \texttt{paper/figures\_src/make\_tables\_exp3.py} from
\texttt{results/exp3\_summary.csv}.
Seeds $3$ and $4$ were added after the registered runs (pre-registration section~11) and
enter the tables (the number of seeds is shown); the registered statements of
Section~\ref{sec:real} keep their three-seed values. The added seeds agree on every direction:
in the low-accessibility regime the informative-context probe gain is $0.026$ and $0.061$
against $0.329$ and $0.462$ for the comparator, reversal accuracy $0.07$ against $0.91$--$0.92$,
the within-seed rate response is $+0.032$ and $+0.054$ at $\gamma=30$ and $+0.086$ and $+0.082$ at
$\gamma=10$, and retraining recovers $0.76$ and $0.62$ from the $\kappa=1/256$ encoders at
$\gamma=10$ against $0.64$ and $0.49$ from $\kappa=1$; in the high-accessibility regime a
retrained head recovers $0.89$--$0.93$ from every encoder.
\IfFileExists{sections_iclr_v2/appendix_exp3_tables.tex}{\begin{table}[htbp]\centering\scriptsize
\caption{Experiment 3, unready (whitened, $\gamma=30$) regime, at $L^\ast=0.30$: mean$\pm$sd over seeds (number of seeds). Probe gain = discriminant accuracy on the encoder output of validation centres minus its initial value; accuracies on validation (val) and test (test) patients under the paired conditions.}\label{tab:exp3_unready}
\begin{tabular}{llrcccccc}\toprule
arm & $M$ / $\kappa$ & step & probe gain & rev.\ (val) & ctx-rand.\ (val) & spec.-only (val) & rev.\ (test) & ctx-rand.\ (test) \\ \midrule
\texttt{sp} & 8 & 126$\pm$51 & 0.015$\pm$0.031 & 0.083$\pm$0.014 & 0.500$\pm$0.012 & 0.485$\pm$0.028 & 0.083$\pm$0.005 & 0.496$\pm$0.000 (3) \\
\texttt{sp} & 32 & 113$\pm$22 & 0.032$\pm$0.042 & 0.072$\pm$0.006 & 0.500$\pm$0.010 & 0.510$\pm$0.006 & 0.070$\pm$0.003 & 0.495$\pm$0.001 (3) \\
\texttt{sp} & 128 & 100$\pm$12 & 0.026$\pm$0.017 & 0.075$\pm$0.003 & 0.500$\pm$0.010 & 0.495$\pm$0.015 & 0.073$\pm$0.006 & 0.496$\pm$0.004 (3) \\
\texttt{sp} & 512 & 69$\pm$8 & -0.012$\pm$0.022 & 0.066$\pm$0.001 & 0.498$\pm$0.011 & 0.486$\pm$0.010 & 0.067$\pm$0.005 & 0.495$\pm$0.004 (3) \\
\midrule
\texttt{headlr} & 32 / 0.00390625 & 185$\pm$43 & 0.046$\pm$0.052 & 0.075$\pm$0.012 & 0.500$\pm$0.013 & 0.512$\pm$0.011 & 0.072$\pm$0.002 & 0.494$\pm$0.001 (3) \\
\texttt{headlr} & 32 / 0.0625 & 177$\pm$40 & 0.044$\pm$0.051 & 0.074$\pm$0.011 & 0.500$\pm$0.013 & 0.512$\pm$0.010 & 0.073$\pm$0.001 & 0.495$\pm$0.002 (3) \\
\texttt{headlr} & 32 & 113$\pm$22 & 0.032$\pm$0.042 & 0.072$\pm$0.006 & 0.500$\pm$0.010 & 0.510$\pm$0.006 & 0.070$\pm$0.003 & 0.495$\pm$0.001 (3) \\
\midrule
\texttt{ctxfree} & 8 & 16780$\pm$2138 & 0.352$\pm$0.008 & 0.925$\pm$0.015 & 0.923$\pm$0.015 & 0.924$\pm$0.010 & 0.935$\pm$0.023 & 0.934$\pm$0.023 (3) \\
\texttt{ctxfree} & 128 & 14954$\pm$488 & 0.336$\pm$0.014 & 0.911$\pm$0.030 & 0.908$\pm$0.029 & 0.926$\pm$0.028 & 0.930$\pm$0.014 & 0.928$\pm$0.009 (3) \\
\midrule
\texttt{frozen} & 32 & 395$\pm$95 & 0.000$\pm$0.000 & 0.076$\pm$0.006 & 0.500$\pm$0.009 & 0.515$\pm$0.008 & 0.076$\pm$0.002 & 0.496$\pm$0.002 (3) \\
\bottomrule
\end{tabular}\end{table}
\begin{table}[htbp]\centering\scriptsize
\caption{Experiment 3, unready (whitened, $\gamma=10$; amendment) regime, at $L^\ast=0.30$: mean$\pm$sd over seeds (number of seeds). Probe gain = discriminant accuracy on the encoder output of validation centres minus its initial value; accuracies on validation (val) and test (test) patients under the paired conditions.}\label{tab:exp3_unready10}
\begin{tabular}{llrcccccc}\toprule
arm & $M$ / $\kappa$ & step & probe gain & rev.\ (val) & ctx-rand.\ (val) & spec.-only (val) & rev.\ (test) & ctx-rand.\ (test) \\ \midrule
\texttt{sp} & 32 & 592$\pm$70 & 0.076$\pm$0.055 & 0.067$\pm$0.004 & 0.501$\pm$0.011 & 0.512$\pm$0.007 & 0.063$\pm$0.002 & 0.496$\pm$0.002 (3) \\
\midrule
\texttt{headlr} & 32 / 0.00390625 & 1579$\pm$318 & 0.134$\pm$0.072 & 0.066$\pm$0.010 & 0.501$\pm$0.011 & 0.508$\pm$0.015 & 0.065$\pm$0.005 & 0.496$\pm$0.001 (3) \\
\texttt{headlr} & 32 / 0.0625 & 1337$\pm$227 & 0.124$\pm$0.068 & 0.065$\pm$0.009 & 0.501$\pm$0.011 & 0.508$\pm$0.013 & 0.065$\pm$0.003 & 0.496$\pm$0.001 (3) \\
\texttt{headlr} & 32 & 592$\pm$70 & 0.076$\pm$0.055 & 0.067$\pm$0.004 & 0.501$\pm$0.011 & 0.512$\pm$0.007 & 0.063$\pm$0.002 & 0.496$\pm$0.002 (3) \\
\midrule
\texttt{ctxfree} & 32 & 16003$\pm$231 & 0.350$\pm$0.012 & 0.916$\pm$0.005 & 0.916$\pm$0.006 & 0.925$\pm$0.009 & 0.923$\pm$0.011 & 0.924$\pm$0.011 (3) \\
\midrule
\texttt{frozen} & 32 & 2006$\pm$310 & 0.000$\pm$0.000 & 0.080$\pm$0.008 & 0.502$\pm$0.011 & 0.511$\pm$0.007 & 0.073$\pm$0.004 & 0.495$\pm$0.001 (3) \\
\bottomrule
\end{tabular}\end{table}
\begin{table}[htbp]\centering\scriptsize
\caption{Experiment 3, ready (standardized) regime, at $L^\ast=0.30$: mean$\pm$sd over seeds (number of seeds). Probe gain = discriminant accuracy on the encoder output of validation centres minus its initial value; accuracies on validation (val) and test (test) patients under the paired conditions.}\label{tab:exp3_ready}
\begin{tabular}{llrcccccc}\toprule
arm & $M$ / $\kappa$ & step & probe gain & rev.\ (val) & ctx-rand.\ (val) & spec.-only (val) & rev.\ (test) & ctx-rand.\ (test) \\ \midrule
\texttt{sp} & 8 & 132$\pm$59 & 0.000$\pm$0.017 & 0.087$\pm$0.004 & 0.502$\pm$0.017 & 0.522$\pm$0.035 & 0.088$\pm$0.009 & 0.499$\pm$0.003 (3) \\
\texttt{sp} & 32 & 118$\pm$22 & 0.002$\pm$0.013 & 0.083$\pm$0.013 & 0.502$\pm$0.013 & 0.526$\pm$0.026 & 0.088$\pm$0.013 & 0.500$\pm$0.003 (3) \\
\texttt{sp} & 128 & 103$\pm$20 & 0.006$\pm$0.016 & 0.102$\pm$0.004 & 0.512$\pm$0.011 & 0.578$\pm$0.060 & 0.102$\pm$0.009 & 0.511$\pm$0.008 (3) \\
\texttt{sp} & 512 & 70$\pm$3 & 0.002$\pm$0.014 & 0.092$\pm$0.012 & 0.508$\pm$0.011 & 0.544$\pm$0.017 & 0.092$\pm$0.021 & 0.504$\pm$0.005 (3) \\
\midrule
\texttt{headlr} & 32 / 0.00390625 & 193$\pm$31 & 0.003$\pm$0.013 & 0.083$\pm$0.010 & 0.503$\pm$0.013 & 0.525$\pm$0.028 & 0.086$\pm$0.008 & 0.501$\pm$0.001 (3) \\
\texttt{headlr} & 32 / 0.0625 & 184$\pm$29 & 0.002$\pm$0.013 & 0.084$\pm$0.011 & 0.503$\pm$0.013 & 0.524$\pm$0.028 & 0.086$\pm$0.008 & 0.501$\pm$0.001 (3) \\
\texttt{headlr} & 32 & 118$\pm$22 & 0.002$\pm$0.013 & 0.083$\pm$0.013 & 0.502$\pm$0.013 & 0.526$\pm$0.026 & 0.088$\pm$0.013 & 0.500$\pm$0.003 (3) \\
\midrule
\texttt{ctxfree} & 8 & 1128$\pm$116 & 0.015$\pm$0.021 & 0.929$\pm$0.002 & 0.923$\pm$0.001 & 0.931$\pm$0.005 & 0.953$\pm$0.008 & 0.947$\pm$0.004 (3) \\
\texttt{ctxfree} & 128 & 841$\pm$112 & 0.012$\pm$0.021 & 0.927$\pm$0.006 & 0.925$\pm$0.006 & 0.929$\pm$0.004 & 0.955$\pm$0.001 & 0.954$\pm$0.000 (3) \\
\midrule
\texttt{frozen} & 32 & 442$\pm$114 & 0.000$\pm$0.000 & 0.091$\pm$0.022 & 0.504$\pm$0.013 & 0.529$\pm$0.021 & 0.093$\pm$0.027 & 0.499$\pm$0.006 (3) \\
\bottomrule
\end{tabular}\end{table}
\paragraph{Runs that did not reach $L^\ast$.}
Every run reached $L^\ast$.
}{\emph{[tables pending]}}

\subsection{Five patient folds (extension, registered before its runs)}
Fold~0 is the registered fold; under pre-registration section~12 the same protocols
(per-fold preprocessing and calibration, all regimes and arms, three seeds, both
retraining procedures) were run on the other four patient-level folds of the
companion's five-fold split, for Experiment~3 and for the natural-context
Experiment~4 of Appendix~\ref{app:exp4}. The registered verdicts are not recomputed;
Tables~\ref{tab:exp3_folds} and~\ref{tab:exp4_folds} give the fold-to-fold variation
of the reported contrasts as mean and range, and any contrast whose sign is not the
same in all five folds is named as fold-dependent.
\IfFileExists{sections_iclr_v2/appendix_folds_tables.tex}{% generated by paper/figures_src/make_tables_folds.py
\begin{table}[htbp]\centering\scriptsize
\caption{Experiment 3 across the five patient folds (fold $0^\ast$ is the registered fold; validation patients; seeds 0--2; values at $L^\ast$, means over widths and seeds). Low-accessibility regime unless stated; retrained = reversal accuracy of the retrained head in the high-accessibility regime.}\label{tab:exp3_folds}
\begin{tabular}{lccccccccc}\toprule
fold & probe gain, inf.\ ctx & probe gain, uninf.\ ctx & reversal, inf. & reversal, uninf. & $\Delta$probe $\kappa\,1\to1/256$ ($\gamma=30$) & same, $\gamma=10$ & reversal, high acc. & retrained, random enc. & retrained, $\kappa=1/256$ \\ \midrule
0$^\ast$ & 0.015 & 0.344 & 0.074 & 0.918 & 0.013 & 0.058 & 0.091 & 0.922 & 0.942 \\
\midrule mean [min, max] & 0.015 [0.02, 0.02] & 0.344 [0.34, 0.34] & 0.074 [0.07, 0.07] & 0.918 [0.92, 0.92] & 0.013 [0.01, 0.01] & 0.058 [0.06, 0.06] & 0.091 [0.09, 0.09] & 0.922 [0.92, 0.92] & 0.942 [0.94, 0.94] \\
\bottomrule\end{tabular}\end{table}
\begin{table}[htbp]\centering\scriptsize
\caption{Experiment 4 (non-denoised copy) across the five patient folds (fold $0^\ast$ registered; validation macro-F1 at budget end; seeds 0--2).}\label{tab:exp4_folds}
\begin{tabular}{lccccccccccc}\toprule
fold & per-pixel gain beyond 16 PCs (MLP) & iid, natural & iid, natural 16 PCs & iid, no bottleneck & probe, natural & probe, shuffled & ctx-random, natural & ctx-random, shuffled & retrained, natural enc. & retrained, shuffled enc. & retrained, random enc. \\ \midrule
0$^\ast$ & 0.066 & 0.732 & 0.657 & 0.744 & 0.672 & 0.670 & 0.234 & 0.680 & 0.681 & 0.688 & 0.494 \\
\midrule mean [min, max] & 0.066 [0.07, 0.07] & 0.732 [0.73, 0.73] & 0.657 [0.66, 0.66] & 0.744 [0.74, 0.74] & 0.672 [0.67, 0.67] & 0.670 [0.67, 0.67] & 0.234 [0.23, 0.23] & 0.680 [0.68, 0.68] & 0.681 [0.68, 0.68] & 0.688 [0.69, 0.69] & 0.494 [0.49, 0.49] \\
\bottomrule\end{tabular}\end{table}
}{}
